\documentclass[12pt]{article}
\usepackage{overcite}
\usepackage{fullpage}
\usepackage{amsmath}
\usepackage{mathtools}
\renewcommand{\familydefault}{\sfdefault}
\renewcommand{\thefootnote}{\textit{\alph{footnote}}}

\def\ket#1{| #1 \rangle}
\def\bra#1{\langle #1 |}
\def\aint#1#2{\bra{#1}\ket{#2}}
\def\bm#1{\mbox{\boldmath $#1$}}
\newcommand{\actocc}{\substack{\textrm{act}{}\\\textrm{occ}}}
\newcommand{\allocc}{\substack{\textrm{all}{}\\\textrm{occ}}}
\newcommand{\actvir}{\substack{\textrm{act}{}\\\textrm{vir}}}
\newcommand{\allvir}{\substack{\textrm{all}{}\\\textrm{vir}}}
\newcommand{\allmos}{\substack{\textrm{all}{}\\\textrm{MOs}}}

\begin{document}

\begin{center}
{\bf \large Derivative Integrals for Correlated Gradients and Hessians} \\
\ \\
T. Daniel Crawford \\
{\em Virginia Tech, Blacksburg, Virginia} \\
29 June 2026
\end{center}

\section{First Derivatives}

\noindent
A computationally simple (but inefficient) form of the first derivative of the correlation energy may be expressed as
\begin{equation}
\partial_x E = \sum_{pq} \gamma_{pq} \partial_x f_{pq} + \sum_{pqrs}
\Gamma_{pqrs} \partial_x \aint{pq}{rs},
\label{Eq:dx_E}
\end{equation}
where $\gamma_{pq}$ and $\Gamma_{pq}$ are anti-symmetric, unrelaxed one- and
two-electron density matrices in the spin-orbital basis.
For a real perturbation, $x$, the derivative of the elements of the Fock matrix are given by
\begin{equation}
\partial_x f_{pq} = f_{pq}^{(x)} + U_{pq}^x f_{pp} + U_{qp}^x f_{qq} 
- \frac{1}{2} \sum_{nm}^{\allocc} S_{nm}^{(x)} A_{pnqm}
+ \sum_c^{\allvir} \sum_m^{\allocc} U_{cm}^x A_{pcqm},
\label{Eq:dx_f_pq}
\end{equation}
where $f_{pq}^{(x)}$ denotes the derivative of the AO-basis Fock operator
transformed into the MO basis (the ``skeleton Fock derivatives''), the
$U_{pq}^x$ are coupled-perturbed Hartree-Fock (CPHF) coefficients, and
\begin{equation}
A_{pqrs} \equiv \aint{pq}{rs} + \aint{ps}{rq}.
\end{equation}
Note that $A_{pqrs}$ has permutational symmetries:
\begin{equation}
A_{pqrs} = A_{psrq} = A_{rqps} = A_{rspq} = A_{qpsr} = A_{qrsp} = A_{spqr} = A_{srqp}.
\end{equation}
Separating $\partial_x f_{pq}$ into MO subspaces and choosing the non-canonical perturbed-orbital conditions, 
\begin{equation}
U_{ij}^x = -\frac{1}{2} S_{ij}^{(x)}
\end{equation}
and
\begin{equation}
U_{ab}^x = -\frac{1}{2} S_{ab}^{(x)}
\end{equation}
yields
\begin{equation}
\partial_x f_{ij} = f_{ij}^{(x)} - \frac{1}{2} S_{ij}^{(x)} (f_{ii} + f_{jj})
- \frac{1}{2} \sum_{nm}^{\allocc} S_{nm}^{(x)} A_{injm}
+ \sum_c^{\allvir} \sum_m^{\allocc} U_{cm}^x A_{icjm},
    \label{Eq:dx_f_ij}
\end{equation}
\begin{equation}
\partial_x f_{ab} = f_{ab}^{(x)} - \frac{1}{2} S_{ab}^{(x)} (f_{aa} + f_{bb})
- \frac{1}{2} \sum_{nm}^{\allocc} S_{nm}^{(x)} A_{anbm}
+ \sum_c^{\allvir} \sum_m^{\allocc} U_{cm}^x A_{acbm},
\end{equation}
and
\begin{equation}
\partial_x f_{ai} = \partial_x f_{ia} = f_{ai}^{(x)} + U_{ai}^x (f_{aa} - f_{ii}) - S_{ai}^{(x)} f_{ii} 
- \frac{1}{2} \sum_{nm}^{\allocc} S_{nm}^{(x)} A_{anim}
+ \sum_c^{\allvir} \sum_m^{\allocc} U_{cm}^x A_{acim} = 0.
\end{equation}
where the last equality hold for Hartree-Fock orbitals by construction.  This
is used to obtain the virtual-occupied block of the CPHF coefficients,
\begin{equation}
\mathbf{G} \mathbf{U}^x = \mathbf{B}
\end{equation}
where
\begin{equation}
G_{ai,cm} = (f_{aa} - f_{ii}) \delta_{ac} \delta_{im} + A_{acim}
\end{equation}
and
\begin{equation}
B_{ai}^x = -f_{ai}^{(x)} + S_{ai}^{(x)} f_{ii} + \frac{1}{2} \sum_{nm}^{\allocc} S_{nm}^{(x)} A_{anim}.
\end{equation}
The derivatives of the antisymmetrized two-electron repulsion integrals are
given by
\begin{equation}
\partial_x \aint{pq}{rs} = \aint{pq}{rs}^{(x)} + \sum_t^{\allmos} \left(
U_{tp}^x \aint{tq}{rs} + U_{tq}^x \aint{pt}{rs} + U_{tr}^x \aint{pq}{ts} + U_{ts}^x \aint{pq}{rt}
\right),
\end{equation}
where $\aint{pq}{rs}^{(x)}$ are the skeleton two-electron integral derivatives.

\vspace{0.2in}
\noindent
For frozen-core gradients, all of the above remains unchanged except that the correlation energy is no longer invariant to
rotations between core and active-occupied MOs. Thus, the non-canonical perturbed-orbital condition for the corresponding
$U_{ij}^x$ cannot be used.  Instead, one must explicitly solve for the core-active CPHF coefficients using
Eq.\eqref{Eq:dx_f_ij} under the assumption of canonical MOs, viz., for $i \in$ core and $j \in$ active-occupied,
\begin{equation}
U_{ij}^x = -(f_{ii} - f_{jj})^{-1} \left[ f_{ij}^{(x)} - S_{ij}^{(x)} f_{jj} - \frac{1}{2} \sum_{nm}^{\allocc}
S_{nm}^{(x)} A_{injm} +
\sum_c^{\allvir} \sum_m^{\allocc} U_{cm}^{x} A_{icjm}\right].
\end{equation}

\section{Second Derivatives}

\noindent
Differentiating Eq.~\eqref{Eq:dx_E}, we obtain
\begin{equation}
\partial_{xy} E = \partial_y \partial_x E = \sum_{pq} \left[ \partial_x (\gamma_{pq}) \partial_y f_{pq} + \gamma_{pq}
\partial_{xy} f_{pq} \right] + \sum_{pqrs} \left[ \partial_x (\Gamma_{pqrs}) \partial_y \aint{pq}{rs} + \Gamma_{pqrs}
\partial_{xy} \aint{pq}{rs} \right],
\end{equation}
where we have used parentheses to clarify that the derivative acts only on the function to its immediate right, i.e.\
we have accounted for the chain rule.  The expressions for the derivatives of the one- and two-electron densities depend
on the choice of quantum chemical method, so we will not discuss them further.  The first derivatives of $f_{pq}$ and
$\aint{pq}{rs}$ are given in the previous section, so we are left only with the second-derivatives.

\vspace{0.2in}
\noindent
For the Fock derivatives, instead of Eq.~\eqref{Eq:dx_f_pq}, we start from a more general form that does not assume
canonical orbitals and does not take advantage of the orthonormality condition, viz.
\begin{equation}
\partial_x f_{pq} = f_{pq}^{(x)} + \sum_r^{\allmos} \left( U_{rp}^x f_{rq} + U_{rq}^x f_{pr} \right) + \sum_r^{\allmos}
\sum_m^{\allocc} U_{rm}^x A_{prqm}.
\end{equation}
Differentiating this expression under the assumption of canonical MOs yields
\begin{align}
\partial_{xy} f_{pq} &= f_{pq}^{(xy)}
+ U_{pq}^{xy} \left(f_{pp} - f_{qq}\right) - \xi_{pq}^{xy} f_{qq}
- \frac{1}{2} \sum_n^{\allocc}\sum_m^{\allocc} \xi_{nm}^{xy} A_{pnqm}
+ \sum_c^{\allvir}\sum_m^{\allocc} U_{cm}^{xy} A_{pcqm}
\nonumber \\
& + \sum_r^{\allmos} \left( U_{rp}^x f_{rq}^{(y)} + U_{rp}^y f_{rq}^{(x)} + U_{rq}^x f_{pr}^{(y)} + U_{rq}^y f_{pr}^{(x)} \right)
+ \sum_{rs}^{\allmos} \left( U_{rp}^x U_{sq}^y + U_{rp}^y U_{sq}^x \right) f_{rs}
\nonumber \\
& + \sum_{rs}^{\allmos}\sum_m^{\allocc} U_{rm}^x U_{sm}^y A_{prqs}
+ \sum_r^{\allmos} \sum_m^{\allocc} \left( U_{rm}^x A_{prqm}^{(y)} + U_{rm}^y A_{prqm}^{(x)} \right)
\nonumber \\
& + \sum_{rs}^{\allmos}\sum_m^{\allocc} \left( U_{rp}^x U_{sm}^y + U_{rp}^y U_{sm}^x \right) A_{rsqm}
+ \sum_{rs}^{\allmos}\sum_m^{\allocc} \left( U_{rq}^x U_{sm}^y + U_{rq}^y U_{sm}^x \right) A_{psrm},
\end{align}
where
\begin{equation}
\xi_{pq}^{xy} \equiv S_{pq}^{xy} + \sum_r^{\allmos} \left( 
U_{qr}^x U_{pr}^y + U_{pr}^x U_{qr}^y - S_{qr}^x S_{pr}^y - S_{pr}^x S_{qr}^y
\right),
\end{equation}
and is obtained from the second derivative of the normalization condition,
\begin{equation}
        U_{pq}^{xy} + U_{qp}^{xy} + \xi_{pq}^{xy} = 0.
\end{equation}
The second derivatives of the two-electron integrals are given by
\begin{align}
\partial_{xy} \aint{pq}{rs} &= \aint{pq}{rs}^{(xy)}
+ \sum_t^{\allmos} \left( U_{tp}^{xy} \aint{tq}{rs} + U_{tq}^{xy} \aint{pt}{rs} + U_{tr}^{xy} \aint{pq}{ts} +
        U_{ts}^{xy} \aint{pq}{rt} \right)
\nonumber \\
& + \sum_t^{\allmos} \left(
U_{tp}^{x} \aint{tq}{rs}^{(y)} + U_{tq}^{x} \aint{pt}{rs}^{(y)} + U_{tr}^{x} \aint{pq}{ts}^{(y)} + U_{ts}^{x}
\aint{pq}{rt}^{(y)}
\right)
\nonumber \\
& + \sum_t^{\allmos} \left(
U_{tp}^{y} \aint{tq}{rs}^{(x)} + U_{tq}^{y} \aint{pt}{rs}^{(x)} + U_{tr}^{y} \aint{pq}{ts}^{(x)} + U_{ts}^{y}
\aint{pq}{rt}^{(x)}
\right)
\nonumber \\
& + \sum_{tu}^{\allmos} \left[
\left( U_{tp}^x U_{uq}^y + U_{tp}^y U_{uq}^x \right) \aint{tu}{rs}
+ \left( U_{tp}^x U_{ur}^y + U_{tp}^y U_{ur}^x \right) \aint{tq}{us}
\right.
\nonumber \\
& \qquad + \left( U_{tp}^x U_{us}^y + U_{tp}^y U_{us}^x \right) \aint{tq}{ru}
+ \left( U_{tq}^x U_{ur}^y + U_{tq}^y U_{ur}^x \right) \aint{pt}{us}
\nonumber \\
& \left.
\qquad + \left( U_{tq}^x U_{us}^y + U_{tq}^y U_{us}^x \right) \aint{pt}{ru}
+ \left( U_{tr}^x U_{us}^y + U_{tr}^y U_{us}^x \right) \aint{pq}{tu}
\right]
\end{align}

\end{document}
